% Generated by `taskgen course examples` from number_theory.wilson (seed 8).
% Do not edit: change examples.yaml and regenerate.
\begin{example}[тест Вильсона, $n = 11$]\label{ex:wilson-test}
С помощью теоремы Вильсона определить, является ли простым число $n = 11$: вычислить $(n - 1)! \bmod n$, приводя по модулю $n$ после каждого умножения. Если $n$ окажется простым, разбить числа $2, 3, \ldots, n - 2$ на пары взаимно обратных по модулю $n$ и объяснить с их помощью полученный ответ.

\par\noindent\textit{Решение.} Теорема Вильсона: $n > 1$ простое $\iff (n - 1)! \equiv -1 \pmod n$. Вычисляем $f_k = k! \bmod 11$ по правилу $f_k = f_{k-1} \cdot k \bmod 11$:\[\begin{array}{c|rrrrrrrrrr}k & 1 & 2 & 3 & 4 & 5 & 6 & 7 & 8 & 9 & 10 \\ \hline k! \bmod 11 & 1 & 2 & 6 & 2 & 10 & 5 & 2 & 5 & 1 & 10\end{array}\]Получили $(11 - 1)! \equiv 10 \equiv -1 \pmod{11}$, т.\,е. $n = 11$ --- простое. Пары $\{a, a^{-1}\}$, $a \ne a^{-1}$: $\{2, 6\}$, $\{3, 4\}$, $\{5, 9\}$, $\{7, 8\}$. Произведение каждой пары $\equiv 1$; из сомножителей $(n - 1)!$ остаются $1$ и $n - 1$ (единственные решения $x^2 \equiv 1$), поэтому $(n - 1)! \equiv 1 \cdot (n - 1) \equiv -1$.
\end{example}
